% Шаблон make_monthly_report. Образец заполнения: reports/monthly/2026-11-12/report.tex. Компиляция: tectonic report.tex
\documentclass[11pt,a4paper]{article}

\usepackage{fontspec}
\usepackage{polyglossia}
\setmainlanguage{russian}
\setotherlanguage{english}
\setmainfont{Times New Roman}
\newfontfamily\cyrillicfont{Times New Roman}
\setsansfont{Helvetica}
\newfontfamily\cyrillicfontsf{Helvetica}
\setmonofont{Menlo}[Scale=0.85]
\newfontfamily\cyrillicfonttt{Menlo}[Scale=0.85]

\usepackage[margin=2cm]{geometry}
\usepackage{amsmath,amssymb}
\usepackage{booktabs,array,multirow,tabularx,longtable}
\usepackage{graphicx}
\usepackage{float}
\usepackage[font=small,labelsep=period]{caption}
\usepackage[shortlabels]{enumitem}
\usepackage{xcolor}
\usepackage{tikz}
\usetikzlibrary{arrows.meta}
\usepackage[colorlinks=true,linkcolor=black,urlcolor=blue!50!black,citecolor=black]{hyperref}

\setlist{nosep,leftmargin=1.2em}
\setlength{\parindent}{0pt}
\setlength{\parskip}{4pt}
\setlength{\emergencystretch}{3em}
\sloppy
\graphicspath{{figures/}}

\definecolor{ml}{HTML}{2A78D6}
\definecolor{ppub}{HTML}{EB6834}

\newcommand{\NMAD}{\mathrm{NMAD}}
\newcommand{\fnref}[1]{\textsuperscript{\ref{#1}}}
\newcolumntype{Y}{>{\raggedright\arraybackslash}X}

\begin{document}

\begin{center}
{\Large Отчёт}\\[4pt]
{\large за период <<13 МЕСЯЦА – 12 МЕСЯЦА ГОД>> г.}\\[8pt]
Кадыров Э.Э., группа 420, кафедра ИИТ, ВМК МГУ\\
Научные руководители: <научные руководители>\\[4pt]
Тема: «Прогноз космологического красного смещения по фотометрическим признакам»
\end{center}

\section{Выполнение календарного плана}
\begin{table}[H]
\centering\small
\begin{tabularx}{\linewidth}{>{\raggedright\arraybackslash}p{2.3cm}>{\raggedright\arraybackslash}p{3.9cm}Y>{\raggedright\arraybackslash}p{1.9cm}}
\toprule
Срок & Задача плана & Что сделано & Статус \\
\midrule
<<СТРОКА НА КАЖДЫЙ ПУНКТ ПРОШЛОГО ПЛАНА>> \\
\bottomrule
\end{tabularx}
\end{table}

\section{Выполненная работа}
% Подразделы по научным вопросам в хронологии месяца, у каждого рисунок (\caption{} пустой).
% Главный результат месяца — подробно: таблица метрик, таблица ΔNMAD с 95% ДИ, опубликованные методы сравнения, слабые места.
\subsection{<<ВОПРОС 1>>}
<<ЧТО СДЕЛАНО: ДАННЫЕ, МОДЕЛЬ, С ЧЕМ СРАВНИВАЛИ>>
\begin{figure}[H]
\centering\includegraphics[width=\linewidth]{<<РИСУНОК>>}
\caption{}
\end{figure}
\begin{itemize}
  \item <<ЧИСЛА С ДИ И ВЫВОД, ВКЛЮЧАЯ ОТРИЦАТЕЛЬНЫЙ>>
\end{itemize}

\subsection{Отчёты, презентация и инструменты}
\begin{itemize}
  \item <<ОТЧЁТЫ С ЧИСЛОМ СТРАНИЦ, ПРЕЗЕНТАЦИЯ СО ССЫЛКОЙ И ЧИСЛОМ СЛАЙДОВ>>
\end{itemize}

\section{Тестовые наборы}
\begin{table}[H]
\centering\small
\begin{tabularx}{\linewidth}{>{\raggedright\arraybackslash}p{3.3cm}lr Y >{\raggedright\arraybackslash}p{3.2cm}}
\toprule
Набор & Поле & $n$ & Метки и отбор & Где использован \\
\midrule
<<ВСЕ НАБОРЫ, НА КОТОРЫХ СЧИТАЛИСЬ МЕТРИКИ>> \\
\bottomrule
\end{tabularx}
\end{table}

\section{Репозиторий и объём кода}
\begin{itemize}
  \item <<ССЫЛКИ (ПРОВЕРЕНЫ), ЧТО ВЫЛОЖЕНО И ЧТО НЕТ>>
  \item <<ЧИСЛО СТРОК ЗА ПЕРИОД ПО КАТЕГОРИЯМ И СПОСОБ ПОДСЧЁТА>>
\end{itemize}

\end{document}
